% Pointwise full-current neighborhood; no standalone preamble.
% Requires amsmath/amssymb/amsthm, theorem environment, and \norm.
% Uses the preceding pointwise baseline theorem. Labels have prefix pc-.
\section{The full flat current in the pointwise weak norms}
\label{pc-section}

We now transfer the pointwise estimates to arbitrary bounded complex
fine currents. The background remains exactly flat. For the corrected
nonlinear constraint term, refinement is $M=L^j$ with fixed odd $L$,
as in the independently audited fixed-step quadratic-composition
theorem. The result does not use a pointwise bound for the projection
$\Pi$ on arbitrary bounded fields and does not differentiate the current.

\subsection{The source norms and spatial weights}

At homogeneous $\ell=1$ and background $U=1$, deposited Raw I,
lines 12740--12752, defines the equivalent norms
\begin{equation}\label{pc-source-norms}
 N(A)=\max\{\norm{A}_\infty,\norm{\nabla A}_\infty\},\qquad
 N_\beta(A)=\max\{N(A),[\nabla A]_\beta\},
 \quad 0\le\beta\le\beta_0<1.
\end{equation}
Here $\nabla$ collects the first fine differences of a vector field;
the H\"older member can be taken on fixed coarse-cell collars.
The pointwise estimates of the preceding section also imply its
global discrete version. There is no Laplacian in $N_\beta$.
The source explicitly excludes the Laplacian member for the current
inverse at Raw I, lines 13759--13765.

Let $\varpi(n)>0$ be a cell weight satisfying
\begin{equation}\label{pc-weight-ratio}
 \frac{\varpi(n)}{\varpi(m)}
       \le e^{\kappa\operatorname{dist}(n,m)}.
\end{equation}
The distance is ordinary or periodic cell distance. Define
$\norm{f}_{\infty,\varpi}$, $N_\varpi$, and $N_{\beta,\varpi}$
by the corresponding weighted supremum over cells and fixed collars.
Neighboring weights have uniformly bounded ratios. In particular
$\varpi(n)=e^{\kappa\operatorname{dist}(n,m_0)}$ is allowed.
These norms are applied to complete fields. We never sharply cut a
differentiable field into cells before applying $D^*$, which would
introduce artificial fine-scale boundary derivatives.

Retain
\begin{equation}\label{pc-scalar-data}
 \mathcal L=D^*D,\qquad Z=G'RG',\qquad
 \Lambda=Z\mathcal L D^*=G'RD^*,\qquad \Pi=I-D\Lambda.
\end{equation}
Stars denote physical adjoints on the real slice and their fixed
analytic bilinear continuations for complex parameters. The
penalized baseline is
$G_0=(d^*d+DRD^*+Q^*aQ)^{-1}$, with matched unit-scale value $a=1$.

\subsection{The scalar smoothing needed for the current}

For a common small $\kappa>0$, the scalar inverse supplies constants
independent of $M$ and compatible volume such that
\begin{align}
 \norm{Z\mathcal L f}_{\infty,\varpi}
       &\le c_0\norm{f}_{\infty,\varpi},&
 \norm{DZ\mathcal L f}_{\infty,\varpi}
       &\le c_1\norm{f}_{\infty,\varpi},
       \label{pc-scalar-bounded}\\
 N_{\beta,\varpi}(DZf)
       &\le z_1\norm{f}_{\infty,\varpi},&
 N_{\beta,\varpi}(DZD^*f)
       &\le z_2\norm{f}_{\infty,\varpi}.
       \label{pc-scalar-holder}
\end{align}
Constants can be chosen at the fixed ceiling $\beta_0$, uniformly
for $0\le\beta\le\beta_0$.

Here is the pointwise proof of these additional scalar facts. The
canceled alias matrix of $Z$ is the high diagonal $P_r^{-2}$ plus
a fixed number of rank terms whose high vectors are $a_r^-/P_r^2$
and their transpose partners. Its central entries have no massless
pole. The high diagonal has the heat representation
\begin{equation}\label{pc-biharmonic-heat}
 \int_0^\infty t\,[H_z(t)-e^{-tP_0}\Pi_0(z)]\,dt.
\end{equation}
The weighted heat derivative bounds extend to every fixed integer
order $k$: split the semigroup into $k$ equal-time factors, put one
fine difference on each, and use weighted $L^1$ convolution.
Thus the small-time bound is $C_kt^{-k/2}$. An additional
$\beta$-H\"older output difference costs $t^{-\beta/2}$ by the
same interpolation used for the baseline inverse. Input derivatives
are derivatives of the kernel, or discrete summation by parts, and
do not require derivatives of the forcing. The relevant integral is
\[
 \int_0^1 t^{1-(m+\beta)/2}\,dt<\infty
                  \quad\text{if }m+\beta<4.
\]
For $Z\mathcal L,DZ\mathcal L$ the derivative orders are two and
three. For the gradient H\"older members of $DZ,DZD^*$ they are
two and three, respectively, before the extra $\beta$. Large times
use the high-alias gap.

For the rank terms, an input row with the required derivatives has
uniform alias $\ell^2$ norm. Output vectors obey
$|a_r/P_r^2|\le Cw(r)/(1+|r|)^4$, where
$w(r)=\prod_\nu(1+|r_\nu|)^{-1}$ and its dyadic shell sum is
$O((j+1)^{d-1})$. Every output derivative/H\"older order strictly
below four is therefore absolutely summable. Physical-cell $L^2$
input rows act on bounded forcing without a dimension factor.
The same periodic holomorphic gluing as for $G_0$ gives exponential
cell kernels and hence \eqref{pc-scalar-bounded}--\eqref{pc-scalar-holder}.
Finite periodic sampling preserves the bounds.

It follows, using the actual derivative of the complete field, that
\begin{equation}\label{pc-lambda-on-N}
 \norm{\Lambda A}_{\infty,\varpi}\le l_0N_\varpi(A),\qquad
 \norm{D\Lambda A}_{\infty,\varpi}\le l_1N_\varpi(A),\qquad
 \norm{\Pi A}_{\infty,\varpi}\le(1+l_1)N_\varpi(A).
\end{equation}
No assertion that $D\Lambda$ or $\Pi$ is bounded on unrestricted
$\ell^\infty$ input is used.

\subsection{Both current terms map $N$ to $N_\beta$}

The matched local auxiliary factorization is
\begin{equation}\label{pc-aux-factorization}
 \Theta'_J=\mathfrak j_J^*\Lambda+\Lambda^*\mathfrak j_J
       +\Lambda^*V_JD\Lambda+(D\Lambda)^*V_J^*\Lambda.
\end{equation}
The finite-incidence coefficient operators $\mathfrak j_J,V_J$
and their analytic transposes have weighted pointwise norms at most
$c_j\norm{J}_\infty,c_v\norm{J}_\infty$. They contain no derivative
of $J$. The exact scalar identity $G_0\Lambda^*=DZ$ gives
\begin{equation}\label{pc-outer-aux-smoothing}
 G_0\Theta'_J
  =G_0\mathfrak j_J^*\Lambda+DZ\mathfrak j_J
       +DZV_JD\Lambda+DZD^*V_J^*\Lambda.
\end{equation}
Let $g_\beta$ bound $G_0$ from weighted bounded forcing to
$N_{\beta,\varpi}$. Then
\begin{equation}\label{pc-aux-N-bound}
 N_{\beta,\varpi}(G_0\Theta'_JA)
        \le C_{\rm aux}\norm{J}_\infty N_\varpi(A),
\end{equation}
where one possible constant is
\[
 C_{\rm aux}=g_\beta c_jl_0+z_1c_j+z_1c_vl_1+z_2c_vl_0.
\]
This is the homogeneous normalized source (J1), including weights.
Its input requires $N$ only, not H\"older regularity of $A$ or $J$.

For the full constraint term, write
\begin{equation}\label{pc-constraint-weight}
 \Theta_J=T_{v_J},\qquad v_J=E_J^T J,
 \qquad
 \langle A,T_vB\rangle_\eta
      =2\langle C_j^{(2)}(A,B),v\rangle_1,
\end{equation}
where $E_J$ is the auxiliary normalized extension and $C_j^{(2)}$
is the quadratic jet of the actual corrected fixed-$L$ composition.
The previously proved auxiliary cell-Schur estimate gives, analytically,
$\norm{v_J}_\infty\le C_v\norm{J}_\infty$ throughout its current ball.

The direct pointwise estimate for $T_v$ uses the same physical
factor as its Hilbert proof. At a step of fine spacing $h$, the
quadratic coefficient is $(h/L)C_L$. Its finite coefficient tensor,
paired with coarse measure and converted to the fine physical
adjoint, has an absolute row sum at most $c_Lh\norm{v}_\infty$.
The factor $L^d$ in this conversion is fixed. The quadratic chain
rule expresses the full operator as
\begin{equation}\label{pc-pointwise-chain}
 T_v=\sum_{t=0}^{j-1}Q_{0:t}^*T_{h_t,v_t}Q_{0:t},\qquad
 h_t=L^{t-j},\qquad v_t=Q_{t+1:j}^*v.
\end{equation}
Both flat $Q$ and its physical adjoint $Q^*$ have nonnegative row
coefficients summing to one. They are $\ell^\infty$ contractions,
including the covering-space multiplicities on short tori. Thus
\[
 \norm{T_v}_{\infty\to\infty}
       \le\frac{c_L}{L-1}\norm{v}_\infty.
\]
The composition has fixed unit-cell range $R_*$, since its reading
radii sum over the same geometric series. Consequently
\begin{equation}\label{pc-T-pointwise}
 \norm{T_vf}_{\infty,\varpi}
       \le\tau_\infty\norm{v}_\infty
                         \norm{f}_{\infty,\varpi}.
\end{equation}
The constant includes $e^{\kappa R_*}$. This is a coefficient-row
argument, not a conversion of a final-cell Hilbert norm into a
pointwise norm.

The identity
\begin{equation}\label{pc-Pi-outer}
 G_0\Pi^*=G_0-DZD^*
\end{equation}
now places every smoothing operator outside the bounded current
coefficients. Combining \eqref{pc-lambda-on-N},
\eqref{pc-T-pointwise}, and \eqref{pc-Pi-outer} gives
\begin{align}
 N_{\beta,\varpi}(G_0\Pi^*\Theta_J\Pi A)
       &\le C_{\rm con}\norm{J}_\infty N_\varpi(A),
       \label{pc-constraint-N-bound}\\
 C_{\rm con}&=(g_\beta+z_2)\tau_\infty C_v(1+l_1).
       \notag
\end{align}
In particular no derivative of $J$ has appeared.

There is also a constrained baseline estimate. Put
\begin{equation}\label{pc-constrained-baseline}
 \Gamma=G_0-H_{\rm src}q_{\rm aux}^{-1}H_{\rm src}^*,
 \qquad q_{\rm aux}=q_M+aI.
\end{equation}
The extension $H_{\rm src}$ is pointwise local into $N_\beta$;
its transpose maps bounded fine fields to bounded coarse fields
by the physical-cell Schur bound; and
$q_{\rm aux}^{-1}=QG_0Q^*$ is a local bounded coarse operator.
Hence $\Gamma$ maps weighted bounded forcing to $N_\beta$ and
\begin{equation}\label{pc-constrained-current}
 N_{\beta,\varpi}(\Gamma\Theta_JA)
      \le C\norm{J}_\infty\norm{A}_{\infty,\varpi}.
\end{equation}

\subsection{A common pointwise current neighborhood}

\begin{theorem}\label{pc-theorem}
Let $\varepsilon_H$ be the established full-current Hilbert/Schur
radius, which is contained in the auxiliary radius. Define
\begin{equation}\label{pc-radius}
 C_*=C_{\rm aux}+C_{\rm con},\qquad
 \varepsilon_{\rm pt}
       =\min\{\varepsilon_H,(2C_*)^{-1}\},
\end{equation}
increasing $C_*$ by fixed norm-convention factors if necessary.
For $\norm{J}_\infty<\varepsilon_{\rm pt}$, put
\begin{equation}\label{pc-full-perturbation}
 \Psi_J=\Theta'_J+\Pi^*\Theta_J\Pi,\qquad
 K_J=G_0\Psi_J.
\end{equation}
The actual full homogeneous flat inverse is the convergent series
\begin{equation}\label{pc-neumann-inverse}
 G_1^Jf=(I-K_J)^{-1}G_0f
                  =\sum_{n\ge0}K_J^nG_0f.
\end{equation}
It is holomorphic from the complex bounded-current ball to operators
from bounded fine forcing into $N_\beta$, uniformly for
$0\le\beta\le\beta_0$. With $g_N$ the baseline bounded-forcing-to-$N$
constant, it obeys
\begin{align}
 N_\varpi(G_1^Jf)
       &\le2g_N\norm{f}_{\infty,\varpi},
       \label{pc-inverse-N}\\
 N_{\beta,\varpi}(G_1^Jf)
       &\le(g_\beta+2C_*\varepsilon_{\rm pt}g_N)
                              \norm{f}_{\infty,\varpi}.
       \label{pc-inverse-Nbeta}
\end{align}
All constants are independent of $M=L^j$ and compatible volume.
For forcing supported in cell $m$, choosing
$\varpi(n)=e^{\kappa\operatorname{dist}(n,m)}$ gives exponential
coarse-cell locality in the pointwise target norms. The auxiliary
inverse is the same result with only $\Theta'_J$ retained.
\end{theorem}

\begin{proof}
Equations \eqref{pc-aux-N-bound} and
\eqref{pc-constraint-N-bound} give
\[
 N_{\beta,\varpi}(K_JA)
          \le C_*\norm{J}_\infty N_\varpi(A).
\]
The target dominates $N_\varpi$, so the ratio is below $1/2$
on \eqref{pc-radius}. The series converges in weighted $N$.
For every term after the zeroth, apply the stronger $N$-to-$N_\beta$
bound to its outside $K_J$ and the $N$ contraction to the remaining
factors. Summing gives \eqref{pc-inverse-N}--\eqref{pc-inverse-Nbeta}.
The coefficient $\Theta_J=T_{E_J^TJ}$ is nonlinear in $J$ but is
analytic on the auxiliary ball with the proved linear norm bound.
Uniform convergence therefore gives Banach-valued holomorphy; no
linear-current differentiation formula is presumed for the full family.

For forcing in one cell, the exponentially weighted solution is in
physical $L^2$ on the infinite lattice. It consequently equals the
independently constructed Hilbert inverse on that input. Equality
for bounded forcing follows by the absolutely summable pointwise
cell kernels and the fixed-$M$ summability of the original operator
products. The finite-volume proof is identical with periodic
distance. The minimum with $\varepsilon_H$ ensures that both
constructions use the same actual source inverse.
\end{proof}

The normalized extensions inherit the same bounds:
$E_1^J=G_1^JQ^*q_1^J$ and $E_J=G_\pi^JQ^*q_J$.
Here $Q^*$ maps bounded coarse fields to bounded fine fields, and
the established coarse Schur bounds make $q_1^J,q_J$ bounded on
weighted coarse $\ell^\infty$. Thus their output $N_\beta$ norms
are uniform on the common ball.

\subsection{What has and has not matched the source}

The input and output norms in \eqref{pc-aux-N-bound} and
\eqref{pc-constraint-N-bound} are the explicit homogeneous flat
$N,N_\beta$ norms, up to fixed equivalent conventions for localized
H\"older seminorms. The current has only a supremum bound. These
estimates and \eqref{pc-inverse-Nbeta} supply the bounded-forcing
members of the current program and their weighted versions.

No current Laplacian estimate is asserted. The source expressly
excludes that member, and the Laplacian of $DZD^*$ is an order-zero
operator whose bounded-input pointwise norm has not been supplied.
The source's separate right-divergence members with H\"older input
are not proved by this section. Such input cannot be sharply
localized by cell indicators; its further analysis needs smooth
physical cutoffs or a direct weighted global H\"older estimate.

The fixed-step corrected-average source identification remains the
one audited from the deposited transcriptions. These results are
homogeneous and flat. They do not establish nonflat source stability,
mixed retained-domain estimates, a changing nonlinear constraint
beyond its exact flat quadratic coefficient, or an interacting
Yang--Mills ultraviolet limit.
